%% bare_jrnl.tex
%% V1.4b
%% 2015/08/26
%% by Michael Shell
%% see http://www.michaelshell.org/
%% for current contact information.
%%
%% This is a skeleton file demonstrating the use of IEEEtran.cls
%% (requires IEEEtran.cls version 1.8b or later) with an IEEE
%% journal paper.
%%
%% Support sites:
%% http://www.michaelshell.org/tex/ieeetran/
%% http://www.ctan.org/pkg/ieeetran
%% and
%% http://www.ieee.org/

%%*************************************************************************
%% Legal Notice:
%% This code is offered as-is without any warranty either expressed or
%% implied; without even the implied warranty of MERCHANTABILITY or
%% FITNESS FOR A PARTICULAR PURPOSE! 
%% User assumes all risk.
%% In no event shall the IEEE or any contributor to this code be liable for
%% any damages or losses, including, but not limited to, incidental,
%% consequential, or any other damages, resulting from the use or misuse
%% of any information contained here.
%%
%% All comments are the opinions of their respective authors and are not
%% necessarily endorsed by the IEEE.
%%
%% This work is distributed under the LaTeX Project Public License (LPPL)
%% ( http://www.latex-project.org/ ) version 1.3, and may be freely used,
%% distributed and modified. A copy of the LPPL, version 1.3, is included
%% in the base LaTeX documentation of all distributions of LaTeX released
%% 2003/12/01 or later.
%% Retain all contribution notices and credits.
%% ** Modified files should be clearly indicated as such, including  **
%% ** renaming them and changing author support contact information. **
%%*************************************************************************


% *** Authors should verify (and, if needed, correct) their LaTeX system  ***
% *** with the testflow diagnostic prior to trusting their LaTeX platform ***
% *** with production work. The IEEE's font choices and paper sizes can   ***
% *** trigger bugs that do not appear when using other class files.       ***                          ***
% The testflow support page is at:
% http://www.michaelshell.org/tex/testflow/



\documentclass[journal]{IEEEtran}
%
% If IEEEtran.cls has not been installed into the LaTeX system files,
% manually specify the path to it like:
% \documentclass[journal]{../sty/IEEEtran}
\usepackage{url}
\usepackage{multirow}
\usepackage{amsmath}
\usepackage{booktabs}
\usepackage{amssymb}
\usepackage{float}
\usepackage{graphicx} 
\usepackage{adjustbox} 
\usepackage{caption}
\usepackage{subcaption}
\usepackage{tabularx}
%\usepackage{xcolor}
\usepackage[table]{xcolor}
\usepackage{comment}
\usepackage{makecell}
% Some very useful LaTeX packages include:
% (uncomment the ones you want to load)
\usepackage{bm}

% *** MISC UTILITY PACKAGES ***
%
%\usepackage{ifpdf}
% Heiko Oberdiek's ifpdf.sty is very useful if you need conditional
% compilation based on whether the output is pdf or dvi.
% usage:
% \ifpdf
%   % pdf code
% \else
%   % dvi code
% \fi
% The latest version of ifpdf.sty can be obtained from:
% http://www.ctan.org/pkg/ifpdf
% Also, note that IEEEtran.cls V1.7 and later provides a builtin
% \ifCLASSINFOpdf conditional that works the same way.
% When switching from latex to pdflatex and vice-versa, the compiler may
% have to be run twice to clear warning/error messages.






% *** CITATION PACKAGES ***
%
%\usepackage{cite}
% cite.sty was written by Donald Arseneau
% V1.6 and later of IEEEtran pre-defines the format of the cite.sty package
% \cite{} output to follow that of the IEEE. Loading the cite package will
% result in citation numbers being automatically sorted and properly
% "compressed/ranged". e.g., [1], [9], [2], [7], [5], [6] without using
% cite.sty will become [1], [2], [5]--[7], [9] using cite.sty. cite.sty's
% \cite will automatically add leading space, if needed. Use cite.sty's
% noadjust option (cite.sty V3.8 and later) if you want to turn this off
% such as if a citation ever needs to be enclosed in parenthesis.
% cite.sty is already installed on most LaTeX systems. Be sure and use
% version 5.0 (2009-03-20) and later if using hyperref.sty.
% The latest version can be obtained at:
% http://www.ctan.org/pkg/cite
% The documentation is contained in the cite.sty file itself.






% *** GRAPHICS RELATED PACKAGES ***
%
\ifCLASSINFOpdf
  % \usepackage[pdftex]{graphicx}
  % declare the path(s) where your graphic files are
  % \graphicspath{{../pdf/}{../jpeg/}}
  % and their extensions so you won't have to specify these with
  % every instance of \includegraphics
  % \DeclareGraphicsExtensions{.pdf,.jpeg,.png}
\else
  % or other class option (dvipsone, dvipdf, if not using dvips). graphicx
  % will default to the driver specified in the system graphics.cfg if no
  % driver is specified.
  % \usepackage[dvips]{graphicx}
  % declare the path(s) where your graphic files are
  % \graphicspath{{../eps/}}
  % and their extensions so you won't have to specify these with
  % every instance of \includegraphics
  % \DeclareGraphicsExtensions{.eps}
\fi
% graphicx was written by David Carlisle and Sebastian Rahtz. It is
% required if you want graphics, photos, etc. graphicx.sty is already
% installed on most LaTeX systems. The latest version and documentation
% can be obtained at: 
% http://www.ctan.org/pkg/graphicx
% Another good source of documentation is "Using Imported Graphics in
% LaTeX2e" by Keith Reckdahl which can be found at:
% http://www.ctan.org/pkg/epslatex
%
% latex, and pdflatex in dvi mode, support graphics in encapsulated
% postscript (.eps) format. pdflatex in pdf mode supports graphics
% in .pdf, .jpeg, .png and .mps (metapost) formats. Users should ensure
% that all non-photo figures use a vector format (.eps, .pdf, .mps) and
% not a bitmapped formats (.jpeg, .png). The IEEE frowns on bitmapped formats
% which can result in "jaggedy"/blurry rendering of lines and letters as
% well as large increases in file sizes.
%
% You can find documentation about the pdfTeX application at:
% http://www.tug.org/applications/pdftex





% *** MATH PACKAGES ***
%
%\usepackage{amsmath}
% A popular package from the American Mathematical Society that provides
% many useful and powerful commands for dealing with mathematics.
%
% Note that the amsmath package sets \interdisplaylinepenalty to 10000
% thus preventing page breaks from occurring within multiline equations. Use:
%\interdisplaylinepenalty=2500
% after loading amsmath to restore such page breaks as IEEEtran.cls normally
% does. amsmath.sty is already installed on most LaTeX systems. The latest
% version and documentation can be obtained at:
% http://www.ctan.org/pkg/amsmath





% *** SPECIALIZED LIST PACKAGES ***
%
%\usepackage{algorithmic}
% algorithmic.sty was written by Peter Williams and Rogerio Brito.
% This package provides an algorithmic environment fo describing algorithms.
% You can use the algorithmic environment in-text or within a figure
% environment to provide for a floating algorithm. Do NOT use the algorithm
% floating environment provided by algorithm.sty (by the same authors) or
% algorithm2e.sty (by Christophe Fiorio) as the IEEE does not use dedicated
% algorithm float types and packages that provide these will not provide
% correct IEEE style captions. The latest version and documentation of
% algorithmic.sty can be obtained at:
% http://www.ctan.org/pkg/algorithms
% Also of interest may be the (relatively newer and more customizable)
% algorithmicx.sty package by Szasz Janos:
% http://www.ctan.org/pkg/algorithmicx




% *** ALIGNMENT PACKAGES ***
%
%\usepackage{array}
% Frank Mittelbach's and David Carlisle's array.sty patches and improves
% the standard LaTeX2e array and tabular environments to provide better
% appearance and additional user controls. As the default LaTeX2e table
% generation code is lacking to the point of almost being broken with
% respect to the quality of the end results, all users are strongly
% advised to use an enhanced (at the very least that provided by array.sty)
% set of table tools. array.sty is already installed on most systems. The
% latest version and documentation can be obtained at:
% http://www.ctan.org/pkg/array


% IEEEtran contains the IEEEeqnarray family of commands that can be used to
% generate multiline equations as well as matrices, tables, etc., of high
% quality.




% *** SUBFIGURE PACKAGES ***
%\ifCLASSOPTIONcompsoc
%  \usepackage[caption=false,font=normalsize,labelfont=sf,textfont=sf]{subfig}
%\else
%  \usepackage[caption=false,font=footnotesize]{subfig}
%\fi
% subfig.sty, written by Steven Douglas Cochran, is the modern replacement
% for subfigure.sty, the latter of which is no longer maintained and is
% incompatible with some LaTeX packages including fixltx2e. However,
% subfig.sty requires and automatically loads Axel Sommerfeldt's caption.sty
% which will override IEEEtran.cls' handling of captions and this will result
% in non-IEEE style figure/table captions. To prevent this problem, be sure
% and invoke subfig.sty's "caption=false" package option (available since
% subfig.sty version 1.3, 2005/06/28) as this is will preserve IEEEtran.cls
% handling of captions.
% Note that the Computer Society format requires a larger sans serif font
% than the serif footnote size font used in traditional IEEE formatting
% and thus the need to invoke different subfig.sty package options depending
% on whether compsoc mode has been enabled.
%
% The latest version and documentation of subfig.sty can be obtained at:
% http://www.ctan.org/pkg/subfig




% *** FLOAT PACKAGES ***
%
%\usepackage{fixltx2e}
% fixltx2e, the successor to the earlier fix2col.sty, was written by
% Frank Mittelbach and David Carlisle. This package corrects a few problems
% in the LaTeX2e kernel, the most notable of which is that in current
% LaTeX2e releases, the ordering of single and double column floats is not
% guaranteed to be preserved. Thus, an unpatched LaTeX2e can allow a
% single column figure to be placed prior to an earlier double column
% figure.
% Be aware that LaTeX2e kernels dated 2015 and later have fixltx2e.sty's
% corrections already built into the system in which case a warning will
% be issued if an attempt is made to load fixltx2e.sty as it is no longer
% needed.
% The latest version and documentation can be found at:
% http://www.ctan.org/pkg/fixltx2e


%\usepackage{stfloats}
% stfloats.sty was written by Sigitas Tolusis. This package gives LaTeX2e
% the ability to do double column floats at the bottom of the page as well
% as the top. (e.g., "\begin{figure*}[!b]" is not normally possible in
% LaTeX2e). It also provides a command:
%\fnbelowfloat
% to enable the placement of footnotes below bottom floats (the standard
% LaTeX2e kernel puts them above bottom floats). This is an invasive package
% which rewrites many portions of the LaTeX2e float routines. It may not work
% with other packages that modify the LaTeX2e float routines. The latest
% version and documentation can be obtained at:
% http://www.ctan.org/pkg/stfloats
% Do not use the stfloats baselinefloat ability as the IEEE does not allow
% \baselineskip to stretch. Authors submitting work to the IEEE should note
% that the IEEE rarely uses double column equations and that authors should try
% to avoid such use. Do not be tempted to use the cuted.sty or midfloat.sty
% packages (also by Sigitas Tolusis) as the IEEE does not format its papers in
% such ways.
% Do not attempt to use stfloats with fixltx2e as they are incompatible.
% Instead, use Morten Hogholm'a dblfloatfix which combines the features
% of both fixltx2e and stfloats:
%
% \usepackage{dblfloatfix}
% The latest version can be found at:
% http://www.ctan.org/pkg/dblfloatfix




%\ifCLASSOPTIONcaptionsoff
%  \usepackage[nomarkers]{endfloat}
% \let\MYoriglatexcaption\caption
% \renewcommand{\caption}[2][\relax]{\MYoriglatexcaption[#2]{#2}}
%\fi
% endfloat.sty was written by James Darrell McCauley, Jeff Goldberg and 
% Axel Sommerfeldt. This package may be useful when used in conjunction with 
% IEEEtran.cls'  captionsoff option. Some IEEE journals/societies require that
% submissions have lists of figures/tables at the end of the paper and that
% figures/tables without any captions are placed on a page by themselves at
% the end of the document. If needed, the draftcls IEEEtran class option or
% \CLASSINPUTbaselinestretch interface can be used to increase the line
% spacing as well. Be sure and use the nomarkers option of endfloat to
% prevent endfloat from "marking" where the figures would have been placed
% in the text. The two hack lines of code above are a slight modification of
% that suggested by in the endfloat docs (section 8.4.1) to ensure that
% the full captions always appear in the list of figures/tables - even if
% the user used the short optional argument of \caption[]{}.
% IEEE papers do not typically make use of \caption[]'s optional argument,
% so this should not be an issue. A similar trick can be used to disable
% captions of packages such as subfig.sty that lack options to turn off
% the subcaptions:
% For subfig.sty:
% \let\MYorigsubfloat\subfloat
% \renewcommand{\subfloat}[2][\relax]{\MYorigsubfloat[]{#2}}
% However, the above trick will not work if both optional arguments of
% the \subfloat command are used. Furthermore, there needs to be a
% description of each subfigure *somewhere* and endfloat does not add
% subfigure captions to its list of figures. Thus, the best approach is to
% avoid the use of subfigure captions (many IEEE journals avoid them anyway)
% and instead reference/explain all the subfigures within the main caption.
% The latest version of endfloat.sty and its documentation can obtained at:
% http://www.ctan.org/pkg/endfloat
%
% The IEEEtran \ifCLASSOPTIONcaptionsoff conditional can also be used
% later in the document, say, to conditionally put the References on a 
% page by themselves.




% *** PDF, URL AND HYPERLINK PACKAGES ***
%
%\usepackage{url}
% url.sty was written by Donald Arseneau. It provides better support for
% handling and breaking URLs. url.sty is already installed on most LaTeX
% systems. The latest version and documentation can be obtained at:
% http://www.ctan.org/pkg/url
% Basically, \url{my_url_here}.




% *** Do not adjust lengths that control margins, column widths, etc. ***
% *** Do not use packages that alter fonts (such as pslatex).         ***
% There should be no need to do such things with IEEEtran.cls V1.6 and later.
% (Unless specifically asked to do so by the journal or conference you plan
% to submit to, of course. )


% correct bad hyphenation here
\hyphenation{op-tical net-works semi-conduc-tor}


\begin{document}
%
% paper title
% Titles are generally capitalized except for words such as a, an, and, as,
% at, but, by, for, in, nor, of, on, or, the, to and up, which are usually
% not capitalized unless they are the first or last word of the title.
% Linebreaks \\ can be used within to get better formatting as desired.
% Do not put math or special symbols in the title.
\title{GCVA: A Multiview Fusion Mechanism for Heterogeneous Data Representations}
%GCVA: 
%Heterogeneous Multiview Leaning for vehicle... IoT Data
%
%
% author names and IEEE memberships
% note positions of commas and nonbreaking spaces ( ~ ) LaTeX will not break
% a structure at a ~ so this keeps an author's name from being broken across
% two lines.
% use \thanks{} to gain access to the first footnote area
% a separate \thanks must be used for each paragraph as LaTeX2e's \thanks
% was not built to handle multiple paragraphs
%

\author{Alberto Durán-López, Daniel Bolaños-Martinez, and Maria Bermudez-Edo% <-this % stops a space
\thanks{A. Durán-López, D. Bolaños-Martinez and M. Bermudez-Edo are with the CITIC-UGR and IATUR institutes, University of Granada, Spain. (e-mail: \{albduranlopez, danibolanos, mbe\}@ugr.es)}% <-this % stops a space
}

% The paper headers
\markboth{IEEE Internet of Things Journal}%
{Durán-López \MakeLowercase{\textit{et al.}}: GCVA: A Multiview Fusion Mechanism for Heterogeneous  Data Representations}
%GCVA: A Multiview Fusion for Uncertainty-Aware IoT Analytics


% The only time the second header will appear is for the odd numbered pages
% after the title page when using the twoside option.
% 
% *** Note that you probably will NOT want to include the author's ***
% *** name in the headers of peer review papers.                   ***
% You can use \ifCLASSOPTIONpeerreview for conditional compilation here if
% you desire.




% If you want to put a publisher's ID mark on the page you can do it like
% this:
%\IEEEpubid{0000--0000/00\$00.00~\copyright~2015 IEEE}
% Remember, if you use this you must call \IEEEpubidadjcol in the second
% column for its text to clear the IEEEpubid mark.



% use for special paper notices
%\IEEEspecialpapernotice{(Invited Paper)}

% make the title area
\maketitle

\begin{abstract}
Multiview learning in Internet of Things environments offers the potential to integrate complementary views to capture complex information dynamics. However, standard fusion mechanisms often fail when specific views exhibit uncertainty, as these mechanisms typically assign fixed importance to all views, disregarding that the information contained in each view may vary dynamically. To address this limitation, we propose Gated Cross-View Attention (GCVA), a view-agnostic fusion mechanism that estimates view-specific uncertainty and dynamically down-weights views with high uncertainty before they corrupt the shared feature space. Then, GCVA integrates the uncertainty estimation with cross-view gated-attention mechanisms to fuse the views. We validate this approach using a real-world License Plate Recognition (LPR) dataset from a smart tourist village, constructing three aligned views: temporal, tabular, and graph, to predict vehicle revisits. We benchmark our proposal against established fusion mechanisms, including Gated Multimodal Units (GMU), attention and hybrid approaches. Results demonstrate that GCVA effectively modulates the influence of high-uncertainty views, compared to baseline mechanisms. Our proposal achieves a Weighted F1-score of 0.730, outperforming state-of-the-art baselines by a margin of 2.8\% to 14.2\%, while reducing training time by 4.7\% to 40.2\%. These findings confirm that explicit uncertainty modeling is important for robust forecasting in real-world IoT scenarios, particularly when the predictive reliability of the views varies.
\end{abstract}


% Note that keywords are not normally used for peerreview papers.
\begin{IEEEkeywords}
Internet of Things (IoT), License Plate Recognition (LPR), Multiview Learning, Heterogeneous Data Fusion
\end{IEEEkeywords}






% For peer review papers, you can put extra information on the cover
% page as needed:
% \ifCLASSOPTIONpeerreview
% \begin{center} \bfseries EDICS Category: 3-BBND \end{center}
% \fi
%
% For peerreview papers, this IEEEtran command inserts a page break and
% creates the second title. It will be ignored for other modes.
\IEEEpeerreviewmaketitle

\section{Introduction}

\IEEEPARstart{D}{ata} fusion across multiple views can be classified into three categories according to the source and structure of the information. $(i)$ Multimodal fusion combines distinct data sources (e.g., audio and video sensors) \cite{lin2024m3sa,ilias2024cross,wu2024wttfnet, zhou2024large}. $(ii)$ Homogeneous multiview integrates the same source but varies attributes in different granularities (e.g., monthly vs. daily aggregation) \cite{chen2025detection, guo2024contrastive, jiang2023multi}. $(iii)$ Heterogeneous multiview integrates a single data source through multiple representations, capturing complementary aspects of the same event (e.g., graph, tabular, and time-series representations) \cite{cao2023multi, yi2022multi}.

Data streams can be effectively organized into a heterogeneous multiview structure \cite{yu2025review}. To illustrate this multiview in smart tourism, consider License Plate Recognition (LPR) systems, where vehicle information can be represented using $(i)$ the time series from the LPRs, capturing traffic dynamics, $(ii)$ tabular features summarizing specific visit details (such as stay duration, entry/exit timestamps, and distance traveled within a touristic area), and $(iii)$ directed weighted graphs that represent movement patterns between camera locations. Each view encodes different inductive biases: time series representations capture short-term demand fluctuations, tabular features contextual metadata, and graph models spatial topology and route structure \cite{yu2025review}. However, these views are not equally informative for every prediction instance. Depending on the context, the information of a specific view may fluctuate, introducing view-level variance that necessitates adaptive mechanisms to dynamically modulate its contributions \cite{hou2024mvqs}. Unlike conventional noise, which typically refers to data corruption or missing data, view-level variance represents the dynamic fluctuation of the information gathered by a view \cite{malawade2022hydrafusion}.


Multiview learning has evolved from concatenation strategies to deep learning fusion \cite{yu2025review}. Gated Multimodal Units (GMU) employ multiplicative gates to adaptively weight each view's contribution \cite{rohanian2019gated,arevalo2020gated}, selectively emphasizing or suppressing views based on relevance. Attention fusion uses learned weights to dynamically combine views \cite{lin2024m3sa,wei2022subattentional}. In multimodal fusion, hybrid architectures combine attention and gating but often allow noisy views to interact before filtering \cite{ilias2024cross}, or derive gating weights directly from learned embeddings, lacking a mechanism to assess view reliability.


To address these limitations, we propose Gated Cross-View Attention (GCVA), a novel fusion mechanism that combines uncertainty-aware gating with cross-view attention. It operates in two steps: first, it assesses the uncertainty of each view based on its information content and calculates its confidence score; second, it uses these scores to weight views in a GMU, preventing noisy views from corrupting the fusion process and combining them with cross-view attention.

We validate this approach on the task of predicting whether a vehicle will return to a tourist area using LPR data. The problem is formulated as a multiview prediction task in which each visit is characterized by $(i)$ a time-series view representing the recent traffic dynamics of the region, $(ii)$ a tabular view encoding contextual attributes, and $(iii)$ a directed graph capturing the driven route between sensors. This predictive capability is relevant for applications such as tourism analytics, parking demand forecasting, and traffic management \cite{duran2026gct,duran2026edrs}. 

The summary of our contributions is:

\begin{itemize}
    \item We propose GCVA, a novel fusion mechanism that combines uncertainty-aware gating with cross-view attention to modulate views prior to the multiview fusion.
    \item We construct a real-world multiview dataset from raw LPR data, transforming it into three heterogeneous views: time-series, tabular data, and graphs.
    \item We demonstrate that GCVA outperforms state-of-the-art fusion mechanisms using the created dataset.
\end{itemize}


\section{Related Work} \label{sec:related_work}

In the view fusion literature, most works focus on multimodal learning, combining distinct physical sources such as text, audio, and video \cite{lin2024m3sa,wei2022subattentional,ilias2024cross,du2022gatedattention, wu2024wttfnet, zhou2024large, narvaez2023minimal}. Recently, multiview learning has emerged as a promising direction, where inputs share the same source but differ in representation. Within this domain, multiview homogeneous approaches, using same-view attributes such as distinct time granularities, are becoming more frequent \cite{chen2025detection, guo2024contrastive, jiang2023multi, shi2022multi}. Conversely, multiview heterogeneous methods, which combine structurally different views (e.g., graphs and tabular data) from the same source, remain less explored \cite{cao2023multi, yi2022multi}.

Early fusion directly concatenates feature representations \cite{guo2024contrastive,jiang2023multi,yi2022multi}, whereas late fusion combines predictions via ensemble methods like voting \cite{shi2022multi,chen2025detection}. These apply to both multimodal and multiview learning. For example, Zhou et al. \cite{zhou2024large} concatenated features to integrate spatial and temporal data for traffic speed prediction. Chen and Fang \cite{chen2025detection} utilized majority voting to aggregate decisions from multiple classifiers for detecting DDoS attacks in IoT networks. While computationally efficient, these approaches assign fixed importance, assuming all views are equally reliable for every sample.

Recent multimodal and multiview learning works focus on adaptive deep fusion to dynamically adjust view weights, primarily through gating or attention. Gating approaches use GMUs to modulate view contributions based on information significance \cite{rohanian2019gated, arevalo2020gated,cao2023multi}. Rodrigues et al. \cite{rodrigues2019multimodal} use gated layers to fuse BERT and CNN embeddings for depression detection. Narváez et al. \cite{narvaez2023minimal} introduced a GMU to mitigate sensor noise, whereas Wu et al. \cite{wu2024wttfnet} used gating to integrate environmental context into trajectory predictions. For attention, Wei et al. \cite{wei2022subattentional} use sub-attentional fusion to capture cross-modal interactions, and Lin et al. \cite{lin2024m3sa} employ channel and cross-attention to weigh multi-scale features for sentiment analysis. These methods either model signal intensity (gating) or view correlation (attention), but rarely both simultaneously.

Recent hybrid architectures, used only in multimodal learning and not in multiview, combine attention and gating sequentially. Ilias et al. \cite{ilias2024cross} propose a pipeline where cross-attention precedes a GMU for view fusion. Du et al. \cite{du2022gatedattention} propose a Gated Attention Fusion Network (GAFN) that applies a gating mechanism after the cross-attention stage. Here, attention first captures cross-view interactions, and subsequently, a gate filters the resulting representation to suppress fine-grained noise. However, in these hybrid solutions, cross-view attention occurs before reliability filtering, allowing less informative views to influence the outcome.


\section{Background} \label{sec:preliminaries}

In this section, we describe the general problem of learning from multiple data representations and explain the scheme we use to make these representations compatible with each other.

\subsection{Problem Formulation}
Let $\mathcal{D} = \{(X_i, y_i)\}_{i=1}^N$ be a dataset of $N$ observations. In many IoT scenarios, a single raw data source contains rich, multifaceted information. For example, a single sensor reading might provide spatial coordinates, a timestamp, and metadata simultaneously. To capture this richness, we decompose each observation $X_i$ into $M$ distinct views, denoted as $X_i = \{x_i^{(1)}, x_i^{(2)}, \dots, x_i^{(M)}\}$.

It is important to note that these views do not necessarily come from different physical devices. A single raw input can be processed to generate multiple perspectives: one view might represent the data as a graph to capture relationships, while another view treats it as a time series to capture trends.

The core challenge lies in effectively integrating these diverse perspectives. Since treating views in isolation ignores their interdependencies, we require a multiview fusion mechanism. Formally, our objective is to construct a composite architecture $\Phi$ capable of encoding and fusing these heterogeneous inputs to produce an accurate prediction $\hat{y}_i$:
\begin{equation}
\hat{y}_i = \Phi(x_i^{(1)}, x_i^{(2)}, \dots, x_i^{(M)}).
\end{equation}
Here, $\Phi$ represents the learnable function that aligns the views and combines their complementary information to approximate the true target $y_i$.


\subsection{Encoding and Projection Scheme} \label{subsec:encoding_scheme}

Processing these diverse views requires a modular strategy assigning each a specialized encoder ($\mathcal{E}_v$). For instance, a Gated Recurrent Unit (GRU) processes temporal sequences, a Multi-Layer Perceptron (MLP) handles static tabular data, and a Graph Neural Network (GNN) models relational structures.

For any view $v$, an encoder $\mathcal{E}_v$ processes the view to extract an embedding vector $\mathbf{z}_v$:
\begin{equation}
    \mathbf{z}_v = \mathcal{E}_v(x_i^{(v)}; \theta_v).
\end{equation}

To handle views with variable structures, such as graphs of arbitrary size without canonical node ordering, the designated encoders must employ permutation-invariant pooling operations (e.g., sum, mean, max pooling or readout functions in GNNs) to aggregate the variable-length features into a fixed-size representation \cite{howxu2019,gilmer2017neural}.

However, a structural challenge arises: distinct encoders typically produce output vectors of varying sizes. For instance, a GNN might output a vector of 64 dimensions, while an MLP outputs a vector of 16 dimensions. To address this dimensionality mismatch, a projection step is used to align all representations into a shared latent space. Because the raw embeddings from the specialized encoders naturally differ in dimension, sharing weights across views is not feasible. Therefore, each view employs an independent, learnable linear transformation $\Pi_v$ (parameterized by view-specific weights) to map its output to a fixed, uniform size $d_f$:
\begin{equation} \label{eq:generate_hv}
    h_v = \Pi_v(\mathbf{z}_v) \in \mathbb{R}^{d_f}.
\end{equation}
This ensures that every view is converted into a standardized vector $h_v$ of the same dimension. The resulting aligned set $H = \{h_1, \dots, h_M\}$ acts as a unified input ready to be processed by a fusion mechanism.


\section{Proposed Method: Gated Cross-View Attention} \label{sec:gcva}


GCVA implements two steps. First, it determines which view is trustworthy by estimating its uncertainty, assigning a confidence score to each view. Second, it allows views to exchange information via attention only if they are considered reliable (i.e., the gate is open). If a view is detected as unreliable, its contribution to the fusion is lowered or suppressed to prevent it from contaminating the reliable views. Figure \ref{fig:gcva} illustrates this method.

\begin{figure*}[!t]
  \centering
  \includegraphics[width=0.96\textwidth]{Images/gcva_updated.png}
  \caption{Proposed GCVA fusion mechanism. The model processes heterogeneous input views by first estimating their individual uncertainty ($\sigma_v^2$) via the Confidence-Aware Gating Mechanism. These estimates generate reliability gates ($g_v$) that dynamically modulate the attention scores, lowering the contribution of noisy views before the final weighted fusion $(z_{\text{fused}})$.}
  \label{fig:gcva}
\end{figure*}

\subsection{Confidence-Aware Gating Mechanism}

Before integrating a view representation $h_v$, the model must assess its level of noise. To achieve this, we employ a probabilistic approach where we estimate the aleatoric uncertainty associated with each view. We use a neural head, $f_{\text{unc}}$, to analyze the data and predict a variance term $\sigma_v^2$. If the predicted variance is high, it implies the view contains noise or ambiguity, whereas a low variance suggests the signal is clear. To ensure stability, we apply the Softplus function, guaranteeing a positive estimated variance. Formally, after projecting the input feature $z_v$ to a shared latent space $h_v \in \mathbb{R}^{d_f}$ (Eq. \ref{eq:generate_hv}), the estimated variance is computed as:
\begin{equation}
\sigma_v^2 = \operatorname{softplus}(f_{\text{unc}}(h_v)) + \epsilon,
\end{equation}
where $f_{\text{unc}}: \mathbb{R}^{d_f} \to \mathbb{R}$ is a linear projection layer and $\operatorname{softplus}(x) = \log(1 + e^x)$. The term $\epsilon$ represents a small stability constant. Unlike epistemic uncertainty, which captures model weight variance and typically requires multiple stochastic passes (e.g., Monte Carlo dropout), aleatoric uncertainty represents inherent data noise \cite{kendall2017uncertainties}. Because it is a direct function of the input data, a single deterministic forward pass is sufficient for the network to output this data-dependent variance $\sigma_v^2$. From this variance, we derive a confidence signal $\gamma_v$. Since variance represents noise, confidence is inversely related to it. We define confidence as the logarithmic inverse of the variance to ensure that views with high noise result in low confidence scores:
\begin{equation}
\gamma_v = \log\left(\frac{1}{\sigma_v^2}\right) = -\log(\sigma_v^2).
\end{equation}

After quantifying confidence, we compute a scalar reliability gate $g_v$. This score considers both the semantic content $h_v$ and the certainty level $\gamma_v$. By concatenating the embedding with its confidence and applying a Sigmoid activation $\sigma$, we restrict the output to $(0, 1)$. This allows $g_v$ to act as a trust probability: values near $1$ indicate full trust, whereas values near $0$ silence unreliable views. We compute this gate as:
\begin{equation} \label{eq:gate}
g_v = \sigma\!\bigl(f_{\text{gate}}([\,h_v\,;\,\gamma_v\,])\bigr) \in (0,1),
\end{equation}
where $f_{\text{gate}}: \mathbb{R}^{d_f+1} \to \mathbb{R}$ is a learnable network determining the system's trust in the view based on its confidence.


\subsection{Gated Attention and Fusion}

With the reliability gates established, we fuse the views. Standard attention mechanisms typically allow information to flow freely between views based solely on similarity. However, this is risky because a noisy view might falsely appear relevant to a clean one, corrupting the feature space. To prevent this, we introduce a reliability-aware modulation mechanism. We modify the standard attention equation by injecting our reliability gates to regulate interaction strength. The core idea is that the information exchange intensity between all pairs of views $i$ and $j$ should be determined by their combined confidence. Mathematically, we multiply the standard attention score by the product of their gates $(g_i \cdot g_j)$. This acts as a continuous scaling factor: connection strength is adjusted so that high confidence maximizes information flow and low confidence suppresses it. Let $H \in \mathbb{R}^{M \times d_f}$ be the matrix of stacked view embeddings. We compute the standard projections (Query, Key, Value) of the attention mechanism: $Q = HW_Q$, $K = HW_K$, and $V = HW_V$. We define the gated attention matrix $A \in \mathbb{R}^{M \times M}$, where each element $A_{ij}$ represents the modulated interaction score between views $i$ and $j$:
\begin{equation}
A_{ij} = (Q_i K_j^\top) \cdot (g_i \cdot g_j).
\end{equation}


The final representation $\widetilde{H}$ is updated via a residual connection, adding the learned interactions to the original knowledge:
\begin{equation}
\widetilde{H} = H + \operatorname{softmax}\left(\frac{A}{\tau}\right)V,
\end{equation}
where $\tau$ regulates the sharpness of the attention distribution. In our architecture, we adopt the standard scaling factor $\tau = \sqrt{d_f}$ \cite{vaswani2017attention} to normalize the dot-product magnitude in high-dimensional latent spaces, preventing the Softmax function from entering saturation regions with vanishing gradients.

The final stage involves synthesizing the refined views into a single representation vector. Instead of a simple average, which would dilute the signal with the noise of weaker views, we reuse our reliability gates to perform a weighted consensus. We convert the gates $g_v$ into normalized weights $\omega_v$ using a Softmax function. This ensures that the most reliable views dominate the final embedding mixture. The fused embedding $z_{\text{fused}}$ is computed as:
\begin{equation}
\omega_v = \frac{\exp(g_v)}{\sum_{u = 1}^{M}} \exp(g_u), \qquad z_{\text{fused}} = \sum_{v=1}^{M} \omega_v \,\widetilde{h}_v.
\end{equation}

Finally, the fused vector $z_{\text{fused}}$ is projected to a scalar probability $\hat{y}$ via a sigmoid classifier, ensuring the final decision relies on a robust consensus of the most reliable views.


\section{Case Study}  \label{sec:case_study}

We validate our proposal through a case study based on real-world data. Specifically, we consider a smart tourism area equipped with four license plate recognition (LPR) cameras and use historical records to predict whether a vehicle newly detected by an LPR will return in subsequent visits to the tourism area. This section details the case study, the view construction and the implementation schema.



\subsection{Dataset} \label{subsec:case_study}

The experimental data originates from the Alpujarra region in southern Spain, specifically covering the municipalities of Pampaneira, Capileira, and Bubion. The infrastructure consists of Hikvision IP cameras equipped with deep learning-based LPRs. Each camera features a 2 MP resolution sensor, varifocal lenses (2.8--12 mm focal length), and infrared LEDs with a 50-meter detection range, enabling reliable vehicle detection across varying perspectives and lighting conditions. As illustrated in Figure \ref{fig:map_as_cameras}, four cameras are distributed across the road network: PAM1 and PAM2 monitor the main access points to Pampaneira, while BUB and CAP cover the municipalities of Bubion and Capileira, respectively.

The data consist of timestamped detections at each camera, including the vehicle license plate string and travel direction. We applied preprocessing steps to correct digit detection errors using Damerau-Levenshtein distance \cite{duran2024routerecoverer}, resulting in a set of 450,759 distinct vehicles. From these records, we construct discrete visits, where a visit is defined as a continuous sequence of detections, with a minimum 24-hour gap required to distinguish subsequent visits. We use the first six months as the input data and the subsequent months as the ground truth verification period, labeling a vehicle that returns as ``Repeater'' and one that does not as ``Non-Repeater''. The dataset used for the experiments has been made publicly available in IEEE DataPort\footnote{\url{https://dx.doi.org/10.21227/77hy-hq94}}.

\begin{figure}[H]
  \centering
  \includegraphics[width=0.85\linewidth]{Images/MAPA_DATAGEN.png}
  \caption{Deployment area in the Alpujarra region.}
  \label{fig:map_as_cameras}
\end{figure}


\subsection{View Construction} \label{subsec:view_instantiation}

For each vehicle visit $i$, we instantiate three views: a temporal sequence capturing traffic trends, a contextual feature vector containing visit attributes, and a relational graph representing the sequence of detections between cameras.



\subsubsection{Temporal View (Sequential)}
We utilize a time series of daily aggregated detections. Let $c_t$ denote the number of unique vehicles detected across the entire network on day $t$. We normalize this global series using z-score optimization as $\tilde{c}_t = \frac{c_t - \mu_{\text{train}}}{\sigma_{\text{train}}}$, where $\mu_{\text{train}}$ and $\sigma_{\text{train}}$ are the mean and standard deviation computed over days in the training partition.

For a specific visit $i$ on date $d_i$, we associate the event with the state of the network leading up to it. Specifically, we extract a sliding window of the past $W$ days:
\begin{equation}
\mathbf{s}_i = \big[\tilde{c}_{d_i - (W-1)}, \dots, \tilde{c}_{d_i}\big] \in \mathbb{R}^{W}.
\end{equation}

This creates a local context sequence for each sample, allowing the model to discern whether the visit occurred during a period of high or low congestion. The window size $W$ enables the model to learn inherent periodicities and short-term seasonal demand, which are critical for spatiotemporal forecasting \cite{luca2021deep}. 

\subsubsection{Tabular View (Contextual)}

To enrich the raw data, we extended the dataset with derived contextual information for each visit. As detailed in Table \ref{tab:features_desc}, these attributes capture mobility metrics and external calendar context, providing a semantic layer that complements the raw spatiotemporal data. To capture static metadata, we construct a dense feature vector $\mathbf{x}_i^{\text{tab}} \in \mathbb{R}^{13}$ using the attributes detailed in Table \ref{tab:features_desc}. This view provides context that is not explicit in the other views.


\begin{table}[H]
    \centering
    \caption{Tabular Features used in the Case Study}
    \label{tab:features_desc}
    \footnotesize 
    \renewcommand{\arraystretch}{1.25} 
    \begin{tabularx}{0.97\linewidth}{@{}l l X@{}} 
    \toprule
    \textbf{Feature Name} & \textbf{Type} & \textbf{Description} \\
    \midrule
    Total Distance & Float & Total distance traveled (km). \\
    Visit Time & Float & Total duration of the visit (hours). \\
    Nights & Integer & Number of nights spent in the region. \\
    Entry Day of Week & Integer & Entry day (0=Monday, 6=Sunday). \\
    Exit Day of Week & Integer & Exit day (0=Monday, 6=Sunday). \\
    Entry is Weekend & Binary & 1 if entry is on weekend, 0 otherwise. \\
    Exit is Weekend & Binary & 1 if exit is on weekend, 0 otherwise. \\
    Entry is Holiday & Binary & 1 if entry is a holiday, 0 otherwise. \\
    Exit is Holiday & Binary & 1 if exit is a holiday, 0 otherwise. \\
    Entry High Season & Binary & 1 if entry is in peak season, 0 otherwise. \\
    Exit High Season & Binary & 1 if exit is in peak season, 0 otherwise. \\
    Entry Hour & Integer & Hour of the day of entry (0--23). \\
    Exit Hour & Integer & Hour of the day of exit (0--23). \\
    \bottomrule
    \end{tabularx}
\end{table}


\subsubsection{Graph View (Relational)}
We represent the spatial route of the vehicle as a directed graph $G_i = (V, E_i, X_v, X_e)$ \cite{duran2025route}.
\begin{itemize}
    \item \textbf{Node Set ($V$):} Corresponds to the four fixed camera locations (PAM1, PAM2, BUB, CAP).
    \item \textbf{Edge Set ($E_i$):} Represents the sequential transitions. A directed edge $(u, v)$ is created if the vehicle moved immediately from camera $u$ to camera $v$.
    \item \textbf{Node Features ($X_v$):} Each node is associated with a feature vector $\mathbf{x}_v$ containing a one-hot encoding of the camera id and a binary flag indicating the direction of travel (entry or exit).
    \item \textbf{Edge Features ($X_e$):} Each edge is annotated with the travel time (in minutes) between consecutive detections. These values are then z-score normalized per visit.
\end{itemize}


\subsection{Encoder Instantiation} \label{sec:encoder_instantiation}

We select neural architectures to instantiate the encoder ($\mathcal{E}_v$) for each view. We evaluate the following candidates:

\begin{itemize}
    \item \textbf{Temporal Encoder ($\mathcal{E}_t$):} To capture trends in the sliding window $\mathbf{s}_i$, we employ Long Short-Term Memory (LSTM), Gated Recurrent Unit (GRU), and a Temporal Transformer.     
    \item \textbf{Tabular Encoder ($\mathcal{E}_b$):} To process the feature vector $\mathbf{x}_i^{\text{tab}}$, we evaluate a standard MLP and the TabTransformer.
    \item \textbf{Graph Encoder ($\mathcal{E}_g$):} Given that our graphs include explicit edge attributes (travel times), we employ edge-aware GNNs capable of modulating message passing with relational data. We evaluate GINe, Edge-GCN (EGCN), and Edge-GAT (EGAT).  
\end{itemize}

Finally, to ensure compatibility for the fusion stage, the output $\mathbf{z}_v$ of each encoder is projected to a common embedding dimension $d_f=128$ using view-specific learnable linear layers $\Pi_v$, yielding the inputs $h_g, h_b, h_t$ for the fusion module.



\subsection{Fusion Baselines} \label{subsec:baselines}

To assess the proposed GCVA mechanism, we benchmark it against established fusion strategies using the same set of pre-trained encoders:

\begin{itemize}
    \item \textbf{Late Fusion} (ensemble) \cite{wang2020makes}: Implements decision-level fusion. Computes the final prediction by averaging the output of each view classifier. This baseline quantifies performance without cross-view feature interaction.
    \item \textbf{Concatenation} \cite{baltruvsaitis2018multimodal}: Implements feature-level fusion. It concatenates the projected embeddings $h_v$ into a single vector. This unified representation is then fed into a shared classifier to generate the final prediction. It assumes all views are equally reliable.
    \item \textbf{Gated Multimodal Unit} \cite{arevalo2020gated}: Utilizes learnable gates to weigh the contribution of each modality dynamically based on input features. It focuses on weighing the importance of the entire view. 
    \item \textbf{Attention} \cite{tsai2019multimodal}: Applies a cross-attention mechanism to capture feature-level relationships between views.
    \item \textbf{Cross-Attention with GMU} \cite{ilias2024cross}: It first uses cross-attention to capture feature-level relationships between views, and subsequently employs a GMU to balance the contribution of each view representation, creating a weighted mixture.
    
    \item \textbf{Gated Attention Fusion Network (GAFN)} \cite{du2022gatedattention}: It first uses cross-attention to capture feature relationships between views, and subsequently employs a gate to independently filter each view representation to suppress noise, without acting as a mixing mechanism.
\end{itemize}

Beyond benchmarking performance, GMU and attention baselines act as an ablation study to validate the individual components of our proposal. Specifically, GMU evaluates the impact of gating mechanisms in the absence of cross-view attention, while Attention fusion isolates the benefits of inter-view modeling without the proposed confidence-aware gating mechanism. By comparing these methods, we can effectively decouple the benefits of capturing cross-view dependencies from the advantages provided by uncertainty-aware gating.

\subsection{Implementation and Training Details}

All models were implemented in PyTorch \cite{paszke2019pytorch} and trained on a single NVIDIA GeForce RTX 3090 GPU. We utilized the AdamW optimizer \cite{loshchilov2018decoupled} with a differential learning rate strategy \cite{howard2018universal}: we applied a lower rate of $10^{-4}$ to the encoders to ensure stable feature extraction, and a higher rate of $10^{-3}$ to the fusion modules to accelerate the convergence of the gating and attention mechanisms. The training process spanned 100 epochs with a batch size of 128, employing an early stopping mechanism with a patience of 10 epochs monitoring the validation F1-score. 

All fusion mechanisms share a standardized configuration. Specifically, all graph, temporal, and tabular encoders were instantiated with 2 hidden layers and a hidden dimension of 128, matching the projection size $d_f$ to maintain consistency in the parameter space. Furthermore, all fusion mechanisms (except Late Fusion) feed their output into a 2-layer MLP classifier with ReLU activation. This ensures that performance differences are solely attributable to the fusion mechanism rather than model capacity. Regarding the GCVA configuration, we set the number of attention heads to $H=4$ and the temperature parameter $\tau = \sqrt{d_f}$ to ensure gradient stability in the softmax distribution, consistent with the standard Scaled Dot-Product Attention mechanism \cite{vaswani2017attention}.

Baseline models are optimized using standard Binary Cross-Entropy (BCE). In contrast, GCVA requires a composite loss function to learn the uncertainty parameters that are absent in the baselines. This objective combines the main prediction task with an auxiliary uncertainty regression term based on \cite{kendall2017uncertainties}. The total loss function is a weighted sum of three components:
\begin{equation} \label{eq:total_loss}
    \mathcal{L} = \mathcal{L}_{\text{main}} + \lambda_1\mathcal{L}_{\text{unc}} + \lambda_2\mathcal{L}_{\text{gate}},
\end{equation}
where $\lambda_1, \lambda_2 \in (0,1)$ are hyperparameters optimized on the validation set to balance the contribution of the auxiliary and regularization terms. $\mathcal{L}_{\text{main}} = \text{BCE}(\hat{y}, y)$ is the primary classification loss from the final fused representation. 
To assess the reliability of each of the $M$ views before fusion, each projected embedding $h_v$ is passed through a shared auxiliary classifier to generate an independent, view-specific prediction $\hat{y}_v$. $\mathcal{L}_{\text{unc}}$ then applies the attenuation mechanism to these independent predictions:
\begin{equation} \label{eq:unc_loss}
    \mathcal{L}_{\text{unc}} = \frac{1}{M} \sum_{v=1}^{M} \left( \frac{\text{BCE}(\hat{y}_v, y)}{\sigma_v^2} + \log(\sigma_v^2) \right).
\end{equation}


This formulation down-weights the error for highly uncertain observations, preventing the model from over-penalizing noisy views and learning to align low confidence with unreliable inputs. Meanwhile, the logarithmic term acts as a regularizer to prevent the model from predicting infinite variance. Finally, $\mathcal{L}_{\text{gate}}$ encourages the reliability gates $g_v$ toward distinct binary decisions (0 or 1), acting as a certainty regularizer:
\begin{equation} \label{eq:gate_loss}
    \mathcal{L}_{\text{gate}} = \frac{1}{M} \sum_{v=1}^{M} g_v(1 - g_v).
\end{equation}

By jointly optimizing these three losses, the model reduces the contribution of unreliable views before the final fusion.

\section{Results and Discussion} \label{sec:results_discuss}


\subsection{Sensitivity Analysis} \label{sec:sensitivity_anal}

To empirically validate the selection of the window size ($W$) for the temporal view, we conducted a sensitivity analysis evaluating $W \in \{7, 14, 21, 28, 35, 42\}$ days. As Figure \ref{fig:sensitivity_temporal} shows, predictive performance improves as the historical context expands from one to four weeks, peaking at $W=28$. Extending the window beyond 28 days (e.g., 35 or 42 days) leads to a performance degradation. Furthermore, computational cost increases with the window size. Therefore, we selected $W=28$ for all subsequent experiments.

\begin{figure}[H]
  \centering
  \includegraphics[width=0.96\linewidth]{Images/temporal_window_sensitivity.png}
  \caption{Sensitivity analysis of the temporal window size.}
  \label{fig:sensitivity_temporal}
\end{figure}

Furthermore, to provide insight into the stability of the training process, we added a sensitivity analysis evaluating the model's Weighted F1-score on the validation set across a uniform grid of values for both loss hyperparameters: $\lambda_1, \lambda_2 \in \{0.00, 0.25, 0.50, 0.75, 1.00\}$. We plotted the results in a heatmap (see Figure \ref{fig:sensitivity_heatmap}), showing that our GCVA model achieves its maximum performance at $\lambda_1=0.50$ and $\lambda_2=0.25$. This confirms that the uncertainty regression term ($\lambda_1$) requires a higher weight to properly model the variance of the data, while the gate regularization term ($\lambda_2$) functions as a softer penalty to prevent premature polarization of the reliability gates during early training epochs. We use these values in the subsequent experiments.

\begin{figure}[H]
  \centering
  \includegraphics[width=0.96\linewidth]{Images/gcva_sensitivity_heatmap_asymmetric.png} 
  \caption{Sensitivity analysis of the GCVA loss hyperparameters ($\lambda_1$ and $\lambda_2$).}
  \label{fig:sensitivity_heatmap}
\end{figure}




\subsection{Single-View and Multiview Fusion Performance} \label{sec:single_and_fusion}

We first evaluate the predictive performance of the view-specific encoders, as detailed in Table \ref{tab:classification_report_individual}. Among the graph-based models, the EGAT model achieved the highest performance with a Weighted F1-score of $0.576$, outperforming the $0.548$ of EGCN and $0.472$ of GINe. In the tabular view, the TabTransformer demonstrated superior capacity with an F1-score of $0.566$, compared to $0.464$ for the MLP. Conversely, temporal models (Transformer, GRU, LSTM) exhibited significant limitations, yielding F1-scores between $0.254$ and $0.359$. These models struggled mainly due to the sparse raw data and the difficulty in capturing long-term patterns from aggregated counts alone, indicating that temporal data alone is unreliable for this task. In contrast, Graph and Tabular models performed significantly better because they rely on explicit spatial topology and rich visit metadata, respectively. Based on these results, we selected the best model for each view (EGAT, TabTransformer, and LSTM) to serve as inputs for the fusion experiments.


We present the performance of the fusion mechanisms for both pairwise and three-view settings in Table \ref{tab:classification_report_fusion} and Figure \ref{fig:architecture}. Strategies such as Late Fusion and Concatenation improved upon single-view, with Concatenation reaching an F1-score of $0.679$ in the three-view setting. However, these techniques are inefficient due to the lack of filtering and high computational cost.


Advanced mechanisms yielded further gains. GMU showed a limitation regarding noise robustness: while effective with the Graph + Tabular pair (F1: $0.712$), its performance dropped to an F1-score of $0.692$ when the noisy Temporal view was added. This degradation indicates that standard gating failed to filter the unreliable signal, suggesting the architecture memorized the noise in the temporal sequences rather than suppressing it. Attention-based Fusion performed similarly (F1: $0.701$, AUROC: $0.740$) but required significantly longer training times ($8392$s). 

Cross-Attention with GMU achieved an F1-score of $0.685$ in the three-view setting, struggling to balance the noisy temporal representations after fusion. GAFN performed better with an F1-score of $0.710$, demonstrating the benefit of post-attention filtering, yet it remained inferior to our proposal. In contrast, GCVA achieved the highest overall effectiveness, with a Tabular + Temporal combination peaking at a Weighted F1-score of $0.722$. In the full three-view setting, GCVA yielded the highest AUROC ($0.752$) and an F1-score of $0.730$, outperforming GMU ($0.692$), Attention ($0.701$), and hybrid baselines. This is attributed to uncertainty-aware gating, which reduces the influence of noisy temporal segments before they reach the attention mechanism.


Regarding efficiency, GCVA is highly cost-effective, reducing training time by $4.7\%$ to $40.2\%$ across all fusion mechanisms. It improves upon adaptive baselines by $4.7\%$ (GMU) and $20.0\%$ (Attention), with more gains against hybrid architectures, reducing time by $27.1\%$ relative to Cross-Attention with GMU and $29.6\%$ relative to GAFN. 

Notably, three-view GCVA execution time ($6715$s) exhibits a higher standard deviation ($\pm 942$s). This variability occurs because explicit uncertainty modeling down-weights noisy temporal views at an early stage, which dynamically triggers the early stopping mechanism in fewer epochs compared to baselines that require more time to process the data. As indicated in Table \ref{tab:early_stopping}, the stopping epoch varies across folds ($\pm 5.2$ epochs). Because GCVA requires denser computations per epoch, this epoch variance translates into the observed deviation in total execution time. Therefore, the higher deviation is a direct consequence of the model's ability to avoid unnecessary computations once the gating logic has optimized the view weights.


\begin{table*}[!t]
    \centering
    \caption{Single Model Classification Performance (5-Fold CV)}
    \label{tab:classification_report_individual}
    
    \setlength{\tabcolsep}{3pt} 
    
    \resizebox{\linewidth}{!}{%
    \begin{tabular}{ll rrr r r rrr}
        \toprule
        \textbf{Model} & \textbf{Class} & \multicolumn{1}{c}{\textbf{Precision}} & \multicolumn{1}{c}{\textbf{Recall}} & \multicolumn{1}{c}{\textbf{F1-score}} & \multicolumn{1}{c}{\textbf{AUROC}} & \multicolumn{1}{c}{\textbf{Time(s)}} & \multicolumn{3}{c}{\textbf{Weighted Avg}} \\
        \cmidrule(lr){8-10}
        & & & & & & & \multicolumn{1}{c}{\textbf{Precision}} & \multicolumn{1}{c}{\textbf{Recall}} & \multicolumn{1}{c}{\textbf{F1-score}} \\
        \midrule
        \multicolumn{10}{l}{\textbf{Graph models}} \\
        \midrule
        \multirow{2}{*}{GINe} & Non-Repeater & $0.742 \pm 0.012$ & $0.324 \pm 0.018$ & $0.451 \pm 0.013$ & \multirow{2}{*}{$0.606 \pm 0.014$} & \multirow{2}{*}{$\mathbf{377.88 \pm 20.87}$} & \multirow{2}{*}{$0.617 \pm 0.012$} & \multirow{2}{*}{$0.483 \pm 0.011$} & \multirow{2}{*}{$0.472 \pm 0.010$} \\
            & Repeater        & $0.379 \pm 0.016$ & $0.786 \pm 0.022$ & $0.512 \pm 0.010$ & & & & & \\
        \multirow{2}{*}{EGCN} & Non-Repeater & $0.752 \pm 0.005$ & $0.450 \pm 0.010$ & $0.563 \pm 0.010$ & \multirow{2}{*}{$0.634 \pm 0.017$} & \multirow{2}{*}{$482.99 \pm 14.99$} & \multirow{2}{*}{$0.633 \pm 0.006$} & \multirow{2}{*}{$0.542 \pm 0.006$} & \multirow{2}{*}{$0.548 \pm 0.008$} \\
            & Repeater        & $0.407 \pm 0.016$ & $0.718 \pm 0.007$ & $0.519 \pm 0.012$ & & & & & \\
        \multirow{2}{*}{EGAT} & Non-Repeater & $0.751 \pm 0.004$ & $0.507 \pm 0.007$ & $0.606 \pm 0.010$ & \multirow{2}{*}{$\mathbf{0.643 \pm 0.009}$} & \multirow{2}{*}{$626.11 \pm 17.45$} & \multirow{2}{*}{$\mathbf{0.638} \pm \mathbf{0.003}$} & \multirow{2}{*}{$\mathbf{0.567 \pm 0.006}$} & \multirow{2}{*}{$\mathbf{0.576 \pm 0.009}$} \\
            & Repeater        & $0.421 \pm 0.004$ & $0.681 \pm 0.007$ & $0.520 \pm 0.007$ & & & & & \\
        \midrule
        \multicolumn{10}{l}{\textbf{Tabular models}} \\
        \midrule
        \multirow{2}{*}{MLP} & Non-Repeater & $0.756 \pm 0.019$ & $0.306 \pm 0.013$ & $0.436 \pm 0.009$ & \multirow{2}{*}{$0.631 \pm 0.007$} & \multirow{2}{*}{$\mathbf{194.96 \pm 4.85}$} & \multirow{2}{*}{$0.626 \pm 0.009$} & \multirow{2}{*}{$0.480 \pm 0.011$} & \multirow{2}{*}{$0.464 \pm 0.007$} \\
            & Repeater        & $0.381 \pm 0.012$ & $0.812 \pm 0.014$ & $0.518 \pm 0.012$ & & & & & \\
        \multirow{2}{*}{TabTransformer} & Non-Repeater & $0.747 \pm 0.004$ & $0.491 \pm 0.013$ & $0.593 \pm 0.010$ & \multirow{2}{*}{$\mathbf{0.633 \pm 0.006}$} & \multirow{2}{*}{$203.36 \pm 4.03$} & \multirow{2}{*}{$\mathbf{0.633 \pm 0.004}$}  & \multirow{2}{*}{$\mathbf{0.558 \pm 0.013}$} & \multirow{2}{*}{$\mathbf{0.566 \pm 0.007}$} \\
            & Repeater        & $0.414 \pm 0.010$ & $0.684 \pm 0.018$ & $0.516 \pm 0.007$ & & & & & \\
        \midrule
        \multicolumn{10}{l}{\textbf{Temporal models}} \\
        \midrule
        \multirow{2}{*}{Transformer} & Non-Repeater & $0.750 \pm 0.040$ & $0.088 \pm 0.030$ & $0.158 \pm 0.025$ & \multirow{2}{*}{$0.580 \pm 0.017$} & \multirow{2}{*}{$3083.79 \pm 224.08$} & \multirow{2}{*}{$0.613 \pm 0.030$} & \multirow{2}{*}{$0.383 \pm 0.026$} & \multirow{2}{*}{$0.280 \pm 0.030$} \\
            & Repeater        & $0.352 \pm 0.030$ & $0.944 \pm 0.033$ & $0.513 \pm 0.040$ & & & & & \\
        \multirow{2}{*}{GRU} & Non-Repeater & $0.763 \pm 0.026$ & $0.063 \pm 0.032$ & $0.117 \pm 0.036$ & \multirow{2}{*}{$0.577 \pm 0.041$} & \multirow{2}{*}{$\mathbf{2165.29 \pm 173.17}$} & \multirow{2}{*}{$\mathbf{0.621 \pm 0.026}$} & \multirow{2}{*}{$0.373 \pm 0.017$} & \multirow{2}{*}{$0.254 \pm 0.024$} \\
            & Repeater        & $0.351 \pm 0.035$ & $0.963 \pm 0.020$ & $0.514 \pm 0.021$ & & & & & \\
        \multirow{2}{*}{LSTM} & Non-Repeater & $0.728 \pm 0.029$ & $0.173 \pm 0.027$ & $0.280 \pm 0.024$ & \multirow{2}{*}{$\mathbf{0.583 \pm 0.023}$} & \multirow{2}{*}{$2959.39 \pm 180.24$} & \multirow{2}{*}{$0.601 \pm 0.022$} & \multirow{2}{*}{$\mathbf{0.416 \pm 0.022}$} & \multirow{2}{*}{$\mathbf{0.359 \pm 0.015}$} \\
            & Repeater        & $0.358 \pm 0.036$ & $0.877 \pm 0.022$ & $0.508 \pm 0.017$ & & & & & \\
        \bottomrule
    \end{tabular}%
    }
\end{table*}


\begin{table*}[!t]
    \centering
    \caption{Fusion Strategies Classification Performance (5-Fold CV). To ensure a fair comparison, all fusion mechanisms utilize the same base encoders: EGAT (Graph), TabTransformer (Tabular), and LSTM (Temporal).}
    \label{tab:classification_report_fusion}
    
    \setlength{\tabcolsep}{2.5pt} 
    
    \resizebox{\linewidth}{!}{%
    \begin{tabular}{ll rrr r r rrr}
        \toprule
        \textbf{Model} & \textbf{Class} & \multicolumn{1}{c}{\textbf{Precision}} & \multicolumn{1}{c}{\textbf{Recall}} & \multicolumn{1}{c}{\textbf{F1-score}} & \multicolumn{1}{c}{\textbf{AUROC}} & \multicolumn{1}{c}{\textbf{Time (s)}} & \multicolumn{3}{c}{\textbf{Weighted Avg}} \\
        \cmidrule(lr){8-10}
        & & & & & & & \multicolumn{1}{c}{\textbf{Precision}} & \multicolumn{1}{c}{\textbf{Recall}} & \multicolumn{1}{c}{\textbf{F1-score}} \\
        \midrule
        
        \multicolumn{10}{l}{\textbf{Late Fusion (ensemble)} \cite{wang2020makes,shi2022multi,chen2025detection}} \\
        \midrule
        \multirow{2}{*}{Graphs + Temporal} & Non-Repeater & $0.747 \pm 0.022$ & $0.520 \pm 0.015$ & $0.613 \pm 0.024$ & \multirow{2}{*}{$0.633 \pm 0.033$} & \multirow{2}{*}{$2790.83 \pm 151.01$} & \multirow{2}{*}{$0.635 \pm 0.017$} & \multirow{2}{*}{$0.570 \pm 0.016$} & \multirow{2}{*}{$0.580 \pm 0.015$} \\
            & Repeater        & $0.421 \pm 0.012$ & $0.665 \pm 0.023$ & $0.516 \pm 0.023$ & & & & & \\
        \multirow{2}{*}{Graphs + Tabular} & Non-Repeater & $0.766 \pm 0.034$ & $0.519 \pm 0.024$ & $0.619 \pm 0.020$ & \multirow{2}{*}{$0.676 \pm 0.017$} & \multirow{2}{*}{$6873.78 \pm 353.70$} & \multirow{2}{*}{$0.652 \pm 0.021$} & \multirow{2}{*}{$0.581 \pm 0.023$} & \multirow{2}{*}{$0.590 \pm 0.016$} \\
            & Repeater        & $0.433 \pm 0.016$ & $0.699 \pm 0.023$ & $0.535 \pm 0.020$ & & & & & \\
        \multirow{2}{*}{Tabular + Temporal} & Non-Repeater & $0.762 \pm 0.023$ & $0.599 \pm 0.018$ & $0.671 \pm 0.016$ & \multirow{2}{*}{$0.672 \pm 0.019$} & \multirow{2}{*}{$3888.71 \pm 84.11$} & \multirow{2}{*}{$0.657 \pm 0.016$} & \multirow{2}{*}{$0.615 \pm 0.013$} & \multirow{2}{*}{$0.624 \pm 0.007$} \\
            & Repeater        & $0.458 \pm 0.025$ & $0.645 \pm 0.031$ & $0.535 \pm 0.027$ & & & & & \\
        \multirow{2}{*}{\shortstack[l]{Graphs + Tabular +\\ Temporal}} & Non-Repeater & $0.780 \pm 0.027$ & $0.608 \pm 0.018$ & $0.683 \pm 0.026$ & \multirow{2}{*}{$0.708 \pm 0.017$} & \multirow{2}{*}{$8888.04 \pm 401.73$} & \multirow{2}{*}{$0.674 \pm 0.015$} & \multirow{2}{*}{$0.630 \pm 0.013$} & \multirow{2}{*}{$0.639 \pm 0.023$} \\
            & Repeater        & $0.474 \pm 0.019$ & $0.673 \pm 0.011$ & $0.556 \pm 0.039$ & & & & & \\
        \midrule
        
        \multicolumn{10}{l}{\textbf{Concatenation} \cite{baltruvsaitis2018multimodal,guo2024contrastive,jiang2023multi,yi2022multi,zhou2024large}} \\
        \midrule
        \multirow{2}{*}{Graphs + Temporal} & Non-Repeater & $0.743 \pm 0.022$ & $0.385 \pm 0.038$ & $0.507 \pm 0.018$ & \multirow{2}{*}{$0.619 \pm 0.021$} & \multirow{2}{*}{$\mathbf{1901.94 \pm 153.85}$} & \multirow{2}{*}{$0.621 \pm 0.021$} & \multirow{2}{*}{$0.510 \pm 0.026$} & \multirow{2}{*}{$0.509 \pm 0.016$} \\
            & Repeater        & $0.390 \pm 0.031$ & $0.746 \pm 0.023$ & $0.512 \pm 0.033$ & & & & & \\
        \multirow{2}{*}{Graphs + Tabular} & Non-Repeater & $0.786 \pm 0.011$ & $0.561 \pm 0.026$ & $0.655 \pm 0.041$ & \multirow{2}{*}{$0.695 \pm 0.018$} & \multirow{2}{*}{$8970.88 \pm 416.54$} & \multirow{2}{*}{$0.673 \pm 0.013$} & \multirow{2}{*}{$0.612 \pm 0.018$} & \multirow{2}{*}{$0.621 \pm 0.026$} \\
            & Repeater        & $0.459 \pm 0.023$ & $0.709 \pm 0.029$ & $0.557 \pm 0.024$ & & & & & \\
        \multirow{2}{*}{Tabular + Temporal} & Non-Repeater & $0.779 \pm 0.025$ & $0.628 \pm 0.023$ & $0.696 \pm 0.048$ & \multirow{2}{*}{$0.706 \pm 0.027$} & \multirow{2}{*}{$4123.62 \pm 275.70$} & \multirow{2}{*}{$0.677 \pm 0.009$} & \multirow{2}{*}{$0.640 \pm 0.011$} & \multirow{2}{*}{$0.649 \pm 0.034$} \\
            & Repeater        & $0.483 \pm 0.024$ & $0.662 \pm 0.018$ & $0.559 \pm 0.022$ & & & & & \\

        \multirow{2}{*}{\shortstack[l]{Graphs + Tabular +\\ Temporal}} & Non-Repeater & $0.776 \pm 0.016$ & $0.707 \pm 0.032$ & $0.740 \pm 0.025$ & \multirow{2}{*}{$0.701 \pm 0.010$} & \multirow{2}{*}{$11227.08 \pm 429.14$} & \multirow{2}{*}{$0.688 \pm 0.012$} & \multirow{2}{*}{$0.674 \pm 0.021$} & \multirow{2}{*}{$0.679 \pm 0.014$} \\
            & Repeater        & $0.523 \pm 0.025$ & $0.611 \pm 0.027$ & $0.564 \pm 0.021$ & & & & & \\
        \midrule
        
        \multicolumn{10}{l}{\textbf{Gated (GMU) Fusion} \cite{arevalo2020gated,rohanian2019gated,cao2023multi,rodrigues2019multimodal,narvaez2023minimal}} \\
        \midrule
        \multirow{2}{*}{Graphs + Temporal} & Non-Repeater & $0.775 \pm 0.016$ & $0.542 \pm 0.020$ & $0.638 \pm 0.018$ & \multirow{2}{*}{$0.681 \pm 0.009$} & \multirow{2}{*}{$2905.36 \pm 238.94$} & \multirow{2}{*}{$0.662 \pm 0.012$} & \multirow{2}{*}{$0.597 \pm 0.012$} & \multirow{2}{*}{$0.606 \pm 0.012$} \\
            & Repeater        & $0.446 \pm 0.008$ & $0.700 \pm 0.008$ & $0.545 \pm 0.013$ & & & & & \\
        \multirow{2}{*}{Graphs + Tabular} & Non-Repeater & $0.772 \pm 0.016$ & $0.803 \pm 0.020$ & $0.787 \pm 0.018$ & \multirow{2}{*}{$0.698 \pm 0.014$} & \multirow{2}{*}{$8501.93 \pm 519.01$} & \multirow{2}{*}{$0.711 \pm 0.011$} & \multirow{2}{*}{$0.715 \pm 0.008$} & \multirow{2}{*}{$0.712 \pm 0.011$} \\
            & Repeater        & $0.594 \pm 0.007$ & $0.548 \pm 0.020$ & $0.570 \pm 0.009$ & & & & & \\
        \multirow{2}{*}{Tabular + Temporal} & Non-Repeater & $0.794 \pm 0.011$ & $0.691 \pm 0.021$ & $0.739 \pm 0.011$ & \multirow{2}{*}{$0.736 \pm 0.011$} & \multirow{2}{*}{$5159.00 \pm 156.77$} & \multirow{2}{*}{$0.702 \pm 0.011$} & \multirow{2}{*}{$0.680 \pm 0.014$} & \multirow{2}{*}{$0.686 \pm 0.010$} \\
            & Repeater        & $0.528 \pm 0.014$ & $0.658 \pm 0.009$ & $0.586 \pm 0.013$ & & & & & \\

        \multirow{2}{*}{\shortstack[l]{Graphs + Tabular +\\ Temporal}} & Non-Repeater & $0.784 \pm 0.014$ & $0.723 \pm 0.009$ & $0.752 \pm 0.025$ & \multirow{2}{*}{$0.711 \pm 0.004$} & \multirow{2}{*}{$7044.45 \pm 222.83$} & \multirow{2}{*}{$0.700 \pm 0.009$} & \multirow{2}{*}{$0.688 \pm 0.009$} & \multirow{2}{*}{$0.692 \pm 0.016$} \\
            & Repeater        & $0.541 \pm 0.011$ & $0.622 \pm 0.017$ & $0.579 \pm 0.024$ & & & & & \\
        \midrule
        
        \multicolumn{10}{l}{\textbf{Attention-based Fusion} \cite{tsai2019multimodal,lin2024m3sa,wei2022subattentional,wu2024wttfnet}} \\
        \midrule
        \multirow{2}{*}{Graphs + Temporal} & Non-Repeater & $0.772 \pm 0.013$ & $0.643 \pm 0.012$ & $0.702 \pm 0.017$ & \multirow{2}{*}{$0.684 \pm 0.010$} & \multirow{2}{*}{$3565.19 \pm 81.64$} & \multirow{2}{*}{$0.673 \pm 0.009$} & \multirow{2}{*}{$0.642 \pm 0.011$} & \multirow{2}{*}{$0.650 \pm 0.019$} \\
            & Repeater        & $0.485 \pm 0.012$ & $0.639 \pm 0.019$ & $0.552 \pm 0.025$ & & & & & \\
        \multirow{2}{*}{Graphs + Tabular} & Non-Repeater & $0.788 \pm 0.014$ & $0.729 \pm 0.019$ & $0.757 \pm 0.016$ & \multirow{2}{*}{$0.725 \pm 0.021$} &\multirow{2}{*}{$\mathbf{3148.52 \pm 147.66}$} & \multirow{2}{*}{$0.706 \pm 0.005$} & \multirow{2}{*}{$0.694 \pm 0.012$} & \multirow{2}{*}{$0.698 \pm 0.006$} \\
            & Repeater        & $0.549 \pm 0.015$ & $0.627 \pm 0.007$ & $0.586 \pm 0.021$ & & & & & \\
        \multirow{2}{*}{Tabular + Temporal} & Non-Repeater & $0.783 \pm 0.009$ & $0.714 \pm 0.022$ & $0.747 \pm 0.008$ & \multirow{2}{*}{$0.727 \pm 0.010$} & \multirow{2}{*}{$\mathbf{2383.52 \pm 150.75}$} & \multirow{2}{*}{$0.698 \pm 0.004$} & \multirow{2}{*}{$0.683 \pm 0.015$} & \multirow{2}{*}{$0.688 \pm 0.009$} \\
            & Repeater        & $0.534 \pm 0.013$ & $0.624 \pm 0.003$ & $0.576 \pm 0.020$ & & & & & \\

        \multirow{2}{*}{\shortstack[l]{Graphs + Tabular +\\ Temporal}} & Non-Repeater & $0.795 \pm 0.015$ & $0.722 \pm 0.008$ & $0.757 \pm 0.016$ & \multirow{2}{*}{$0.740 \pm 0.020$} & \multirow{2}{*}{$8392.58 \pm 311.15$} & \multirow{2}{*}{$0.711 \pm 0.012$} & \multirow{2}{*}{$0.696 \pm 0.006$} & \multirow{2}{*}{$0.701 \pm 0.006$} \\
            & Repeater        & $0.550 \pm 0.010$ & $0.647 \pm 0.012$ & $0.595 \pm 0.013$ & & & & & \\
        \midrule


                \multicolumn{10}{l}{\textbf{Cross-Attention with GMU} \cite{ilias2024cross}} \\
        \midrule
        \multirow{2}{*}{Graphs + Temporal} & Non-Repeater & $0.758 \pm 0.019$ & $0.620 \pm 0.018$ & $0.682 \pm 0.021$ & \multirow{2}{*}{$0.665 \pm 0.028$} & \multirow{2}{*}{$7150.20 \pm 237.60$} & \multirow{2}{*}{$0.657 \pm 0.013$} & \multirow{2}{*}{$0.618 \pm 0.012$} & \multirow{2}{*}{$0.630 \pm 0.014$} \\
            & Repeater        & $0.465 \pm 0.028$ & $0.615 \pm 0.009$ & $0.530 \pm 0.022$ & & & & & \\
        \multirow{2}{*}{Graphs + Tabular} & Non-Repeater & $0.775 \pm 0.009$ & $0.740 \pm 0.009$ & $0.757 \pm 0.014$ & \multirow{2}{*}{$0.715 \pm 0.031$} & \multirow{2}{*}{$6920.10 \pm 346.14$} & \multirow{2}{*}{$0.701 \pm 0.013$} & \multirow{2}{*}{$0.690 \pm 0.008$} & \multirow{2}{*}{$0.695 \pm 0.012$} \\
            & Repeater        & $0.560 \pm 0.022$ & $0.595 \pm 0.012$ & $0.577 \pm 0.016$ & & & & & \\
        \multirow{2}{*}{Tabular + Temporal} & Non-Repeater & $0.770 \pm 0.010$ & $0.695 \pm 0.025$ & $0.730 \pm 0.026$ & \multirow{2}{*}{$0.710 \pm 0.028$} & \multirow{2}{*}{$5980.45 \pm 341.06$} & \multirow{2}{*}{$0.680 \pm 0.006$} & \multirow{2}{*}{$0.664 \pm 0.017$} & \multirow{2}{*}{$0.669 \pm 0.017$} \\
            & Repeater        & $0.510 \pm 0.012$ & $0.605 \pm 0.012$ & $0.553 \pm 0.016$ & & & & & \\
        \multirow{2}{*}{\shortstack[l]{Graphs + Tabular +\\ Temporal}} & Non-Repeater & $0.782 \pm 0.018$ & $0.712 \pm 0.026$ & $0.745 \pm 0.013$ & \multirow{2}{*}{$0.715 \pm 0.026$} & \multirow{2}{*}{$9210.55 \pm 368.90$} & \multirow{2}{*}{$0.697 \pm 0.017$} & \multirow{2}{*}{$0.677 \pm 0.016$} & \multirow{2}{*}{$0.685 \pm 0.014$} \\
            & Repeater        & $0.535 \pm 0.021$ & $0.610 \pm 0.011$ & $0.570 \pm 0.024$ & & & & & \\
        \midrule


        \multicolumn{10}{l}{\textbf{GAFN} \cite{du2022gatedattention}} \\
        \midrule
        \multirow{2}{*}{Graphs + Temporal} & Non-Repeater & $0.770 \pm 0.017$ & $0.660 \pm 0.022$ & $0.711 \pm 0.015$ & \multirow{2}{*}{$0.692 \pm 0.024$} & \multirow{2}{*}{$7420.30 \pm 256.95$} & \multirow{2}{*}{$0.674 \pm 0.012$} & \multirow{2}{*}{$0.651 \pm 0.015$} & \multirow{2}{*}{$0.657 \pm 0.014$} \\
            & Repeater        & $0.492 \pm 0.010$ & $0.635 \pm 0.008$ & $0.554 \pm 0.022$ & & & & & \\
        \multirow{2}{*}{Graphs + Tabular} & Non-Repeater & $0.785 \pm 0.012$ & $0.760 \pm 0.029$ & $0.772 \pm 0.020$ & \multirow{2}{*}{$0.732 \pm 0.012$} & \multirow{2}{*}{$7180.80 \pm 520.73$} & \multirow{2}{*}{$0.713 \pm 0.012$} & \multirow{2}{*}{$0.705 \pm 0.021$} & \multirow{2}{*}{$0.708 \pm 0.009$} \\
            & Repeater        & $0.575 \pm 0.023$ & $0.600 \pm 0.008$ & $0.587 \pm 0.019$ & & & & & \\
        \multirow{2}{*}{Tabular + Temporal} & Non-Repeater & $0.788 \pm 0.012$ & $0.755 \pm 0.011$ & $0.771 \pm 0.013$ & \multirow{2}{*}{$0.735 \pm 0.006$} & \multirow{2}{*}{$6250.15 \pm 186.31$} & \multirow{2}{*}{$0.711 \pm 0.012$} & \multirow{2}{*}{$0.705 \pm 0.009$} & \multirow{2}{*}{$0.707 \pm 0.010$} \\
            & Repeater        & $0.565 \pm 0.016$ & $0.610 \pm 0.007$ & $0.586 \pm 0.026$ & & & & & \\
        \multirow{2}{*}{\shortstack[l]{Graphs + Tabular +\\ Temporal}} & Non-Repeater & $0.800 \pm 0.021$ & $0.740 \pm 0.018$ & $0.769 \pm 0.019$ & \multirow{2}{*}{$0.738 \pm 0.011$} & \multirow{2}{*}{$9540.90 \pm 337.86$} & \multirow{2}{*}{$0.717 \pm 0.017$} & \multirow{2}{*}{$0.707 \pm 0.016$} & \multirow{2}{*}{$0.710 \pm 0.012$} \\
            & Repeater        & $0.558 \pm 0.012$ & $0.645 \pm 0.025$ & $0.598 \pm 0.015$ & & & & & \\
        \midrule


        
        \multicolumn{10}{l}{\textbf{GCVA (Ours)}} \\
        \midrule
        \multirow{2}{*}{Graphs + Temporal} & Non-Repeater & $0.775 \pm 0.020$ & $0.680 \pm 0.011$ & $0.725 \pm 0.011$ & \multirow{2}{*}{$\mathbf{0.701 \pm 0.013}$} & \multirow{2}{*}{$7212.92 \pm 330.88$} & \multirow{2}{*}{$\mathbf{0.682 \pm 0.013}$} & \multirow{2}{*}{$\mathbf{0.661 \pm 0.012}$} & \multirow{2}{*}{$\mathbf{0.668 \pm 0.010}$} \\
            & Repeater        & $0.506 \pm 0.015$ & $0.624 \pm 0.016$ & $0.559 \pm 0.015$ & & & & & \\
        \midrule
        \multirow{2}{*}{Graphs + Tabular} & Non-Repeater & $0.790 \pm 0.016$ & $0.772 \pm 0.020$ & $0.781 \pm 0.015$ & \multirow{2}{*}{$\mathbf{0.742 \pm 0.023}$} & \multirow{2}{*}{$7799.87 \pm 68.07$} & \multirow{2}{*}{$\mathbf{0.719 \pm 0.012}$} & \multirow{2}{*}{$\mathbf{0.716 \pm 0.011}$} & \multirow{2}{*}{$\mathbf{0.718 \pm 0.009}$} \\
            & Repeater        & $0.584 \pm 0.015$ & $0.610 \pm 0.016$ & $0.597 \pm 0.017$ & & & & & \\
        \midrule
        \multirow{2}{*}{Tabular + Temporal} & Non-Repeater & $0.794 \pm 0.015$ & $0.775 \pm 0.008$ & $0.784 \pm 0.016$ & \multirow{2}{*}{$\mathbf{0.745 \pm 0.020}$} & \multirow{2}{*}{$6628.50 \pm 100.37$} & \multirow{2}{*}{$\mathbf{0.724 \pm 0.008}$} & \multirow{2}{*}{$\mathbf{0.721 \pm 0.009}$} & \multirow{2}{*}{$\mathbf{0.722 \pm 0.012}$} \\
            & Repeater        & $0.591 \pm 0.015$ & $0.619 \pm 0.017$ & $0.605 \pm 0.010$ & & & & & \\
        \midrule
        \multirow{2}{*}{\shortstack[l]{Graphs + Tabular +\\ Temporal}} & Non-Repeater & $0.826 \pm 0.011$ & $0.754 \pm 0.008$ & $0.788 \pm 0.009$ & \multirow{2}{*}{$\mathbf{0.752 \pm 0.018}$} & \multirow{2}{*}{$\mathbf{6715.40 \pm 942.15}$} & \multirow{2}{*}{$\mathbf{0.740 \pm 0.005}$} & \multirow{2}{*}{$\mathbf{0.726 \pm 0.012}$} & \multirow{2}{*}{$\mathbf{0.730 \pm 0.007}$} \\
            & Repeater        & $0.575 \pm 0.009$ & $0.673 \pm 0.018$ & $0.621 \pm 0.004$ & & & & & \\
        \bottomrule
    \end{tabular}}
\end{table*}


\begin{table}[H]
    \centering
    \caption{Average early stopping epoch across 5-fold CV in the three-view setting. The reported values include a patience of 10 epochs, with a maximum limit of 100 epochs.}
    \label{tab:early_stopping}
    \resizebox{0.84\columnwidth}{!}{%
    \begin{tabular}{l c}
        \toprule
        \textbf{Fusion Mechanism} & \textbf{Early Stopping Epoch} \\
        \midrule
        Late Fusion & $56.1 \pm 2.5$ \\
        Concatenation & $73.9 \pm 2.8$ \\
        Gated (GMU) & $43.4 \pm 1.4$ \\
        Attention-based & $50.0 \pm 1.9$ \\
        Cross-Attention with GMU & $53.4 \pm 2.1$ \\
        GAFN & $54.8 \pm 1.9$ \\
        GCVA (Ours) & $\mathbf{37.2 \pm 5.2}$ \\
        \bottomrule
    \end{tabular}%
    }
\end{table}


\begin{figure*}[!t]
  \centering
  \includegraphics[width=0.96\linewidth]{Images/final_plot_vertical_aligned.png}
  \caption{Weighted F1-score distributions for different fusion strategies and view combinations using 5-fold cross-validation. Boxes represent the interquartile range (IQR), while overlaid points visualize fold-specific results and model stability.}
  \label{fig:architecture}
\end{figure*}


\subsection{Ablation Study} \label{sec:ablation_test}

To evaluate the interaction between semantic features and the confidence score ($g_v = \sigma(f_{\text{gate}}([h_v; \gamma_v]))$), we conducted an ablation study isolating their contributions:

\begin{itemize}
    \item \textbf{GCVA without Confidence ($\gamma_v$):} Formulated as $g_v=\sigma(f_{\text{gate}}(h_v))$, testing if pure semantic representations can independently detect and filter view-specific noise.
    \item \textbf{GCVA without Semantics ($h_v$):} Formulated as $g_v=\sigma(f_{\text{gate}}(\gamma_v))$, evaluating the isolated capability of uncertainty estimation to modulate gates without contextual features.
\end{itemize}

As shown in Table \ref{tab:ablation_gating}, the full GCVA model achieves the highest performance (F1-score: $0.730$, AUROC: $0.752$). Removing the confidence score ($\gamma_v$) degrades the F1-score to $0.712$, and removing semantic features ($h_v$) causes a more severe drop (F1-score: $0.698$). These findings confirm that, although the semantic features contribute more to the final F1-score, both components are complementary.


\begin{table*}[!t]
    \centering
    \caption{Ablation study on the GCVA gating mechanism (Three-Views setting).}
    \label{tab:ablation_gating}
    
    \setlength{\tabcolsep}{2.5pt} 
    
    \resizebox{\linewidth}{!}{%
    \begin{tabular}{ll rrr r r rrr}
        \toprule
        \textbf{Ablation Variant} & \textbf{Class} & \multicolumn{1}{c}{\textbf{Precision}} & \multicolumn{1}{c}{\textbf{Recall}} & \multicolumn{1}{c}{\textbf{F1-score}} & \multicolumn{1}{c}{\textbf{AUROC}} & \multicolumn{1}{c}{\textbf{Time (s)}} & \multicolumn{3}{c}{\textbf{Weighted Avg}} \\
        \cmidrule(lr){8-10}
        & & & & & & & \multicolumn{1}{c}{\textbf{Precision}} & \multicolumn{1}{c}{\textbf{Recall}} & \multicolumn{1}{c}{\textbf{F1-score}} \\
        \midrule
        
        \multirow{2}{*}{GCVA (Full Model)} & Non-Repeater & $0.826 \pm 0.011$ & $0.754 \pm 0.008$ & $0.788 \pm 0.009$ & \multirow{2}{*}{$\mathbf{0.752 \pm 0.018}$} & \multirow{2}{*}{$6715.40 \pm 942.15$} & \multirow{2}{*}{$\mathbf{0.740 \pm 0.005}$} & \multirow{2}{*}{$\mathbf{0.726 \pm 0.012}$} & \multirow{2}{*}{$\mathbf{0.730 \pm 0.007}$} \\
            & Repeater & $0.575 \pm 0.009$ & $0.673 \pm 0.018$ & $0.621 \pm 0.004$ & & & & & \\
        \midrule
        
        \multirow{2}{*}{GCVA w/o Confidence ($\gamma_v$)} & Non-Repeater & $0.811 \pm 0.014$ & $0.735 \pm 0.012$ & $0.771 \pm 0.013$ & \multirow{2}{*}{$0.741 \pm 0.015$} & \multirow{2}{*}{$6512.30 \pm 210.45$} & \multirow{2}{*}{$0.722 \pm 0.011$} & \multirow{2}{*}{$0.708 \pm 0.015$} & \multirow{2}{*}{$0.712 \pm 0.010$} \\
            & Repeater & $0.551 \pm 0.012$ & $0.655 \pm 0.020$ & $0.598 \pm 0.015$ & & & & & \\
        \midrule
        
        \multirow{2}{*}{GCVA w/o Semantics ($h_v$)} & Non-Repeater & $0.802 \pm 0.015$ & $0.720 \pm 0.018$ & $0.759 \pm 0.016$ & \multirow{2}{*}{$0.728 \pm 0.012$} & \multirow{2}{*}{$\mathbf{6425.10 \pm 180.20}$} & \multirow{2}{*}{$0.708 \pm 0.014$} & \multirow{2}{*}{$0.695 \pm 0.012$} & \multirow{2}{*}{$0.698 \pm 0.009$} \\
            & Repeater & $0.540 \pm 0.018$ & $0.638 \pm 0.022$ & $0.584 \pm 0.019$ & & & & & \\
        \bottomrule
    \end{tabular}%
    }
\end{table*}



\subsection{Validation of the Gating Mechanism} \label{sec:validation_gating}

To validate the gating mechanism, we analyze the relationship between learned uncertainty, gate values ($g_v$), and prediction errors. As shown in Fig. \ref{fig:learned_uncertainty}, the model assigns significantly higher uncertainty to the Temporal view ($\mu: 2.81$) compared to the Graph and Tabular views. This aligns with its lower individual performance, as the gating mechanism correctly identifies it as a noise source of information. Fig. \ref{fig:gate_correlation} plots the relationship between the prediction error and the learned gate value ($g_v$) for each view. As observed, there is a negative Pearson correlation across all representations. Specifically, the Temporal view exhibits the strongest correlation ($r = -0.74, p < 0.001$). This confirms that the GCVA architecture dynamically reduces the contribution (lower $g_v$) of a view when its prediction error increases. This behavior ensures that the final fusion is dominated by the most reliable representations at any given time.

To validate the gating mechanism, we analyze the relationship between learned uncertainty, gate values ($g_v$), and prediction errors. Fig. \ref{fig:gate_correlation} plots the relationship between the prediction error and the learned gate value ($g_v$) for each view. As observed, there is a negative Pearson correlation across all representations. Specifically, the Temporal view exhibits the strongest correlation ($r = -0.74, p < 0.001$). This confirms that the GCVA architecture dynamically reduces the contribution (lower $g_v$) of a view when its prediction error increases. As shown in Fig. \ref{fig:learned_uncertainty}, the model assigns significantly higher uncertainty to the Temporal view ($\mu: 2.81$) compared to the Graph and Tabular views. This aligns with its lower individual performance, as the gating mechanism correctly identifies it as a noise source of information. This behavior ensures that the final fusion is dominated by the most reliable representations at any given time.

\begin{figure*}[!t]
  \centering
  \includegraphics[width=0.97\linewidth]{Images/gcva_gate_correlation_new.png}
  \caption{Correlation analysis between the learned gate values ($g_v$) and the view prediction errors.}
  \label{fig:gate_correlation}
\end{figure*}

\begin{figure}[htbp]
  \centering
  \includegraphics[width=0.97\linewidth]{Images/gcva_uncertainty_boxplot.png}
  \caption{Uncertainty distribution across views. The Temporal view shows the highest uncertainty, indicating lower reliability.}
  \label{fig:learned_uncertainty}
\end{figure}




To provide a qualitative case study, Fig. \ref{fig:case_study} shows the hourly aggregated fusion weights ($\omega_v$) over a continuous one-week period. As observed, the GCVA mechanism successfully suppresses the noisy Temporal view ($\omega_{\text{temp}}$), keeping its assigned weight consistently low. The model relies on the Tabular and Graph views, alternating between them depending on their reliability at each hour.

\begin{figure}[htbp] 
    \centering 
    \includegraphics[width=0.96\linewidth]{Images/gcva_1week_case_study_lines.png} 
    \caption{Weekly evolution of view-specific fusion weights ($\omega_v$).} 
    \label{fig:case_study} 
\end{figure}

\subsection{Future Work}
While GCVA filters noisy views by modeling aleatoric uncertainty in a single forward pass, it currently does not account for epistemic uncertainty. We prioritized modeling aleatoric uncertainty because the primary fusion challenges stem from data-dependent noise (e.g., missed sensor detections or varying travel times due to traffic). However, epistemic uncertainty is also relevant in IoT scenarios for detecting rare events or out-of-distribution samples, such as a unique, previously unseen vehicle route. Traditional methods for capturing epistemic uncertainty, such as Monte Carlo sampling or deep ensembles, produce superior predictions but are computationally expensive due to multiple forward passes \cite{ng2022estimating}. Recently, a promising line of research has focused on reducing this computational burden through efficient joint learning at inference \cite{li2023towards} and novel approaches that estimate epistemic uncertainty using a single deterministic network \cite{van2020uncertainty, zanger2025contextual}. In future work, we aim to extend our framework by integrating these cost-effective epistemic solutions, enabling the system to weigh the reliability of views based on both inherent data noise and model uncertainty, without sacrificing inference speed.


\section{Conclusion} \label{sec:conclusion}

In this work, we present GCVA, a novel fusion mechanism designed to prevent noisy views from corrupting the shared feature space. This method operates in two distinct stages: first, it employs a confidence-aware gating mechanism to estimate the uncertainty of each view; second, it utilizes these reliability scores to regulate a gated attention process that modulates information flow. We validated this approach in a smart village case study, demonstrating that relying on single views is insufficient for accurate return visit prediction given the unreliable temporal sequences. Through a comparative analysis of state-of-the-art fusion strategies, we showed that while traditional methods provide marginal gains, they often fail to filter noise from weaker inputs effectively. GCVA overcomes these limitations by dynamically reducing unreliable inputs, achieving F1-score improvements between $2.8\%$ and $14.2\%$ compared to established baselines. Furthermore, it proved highly efficient, reducing total training time by $4.7\%$ to $40.2\%$ across all fusion mechanisms. These findings confirm that explicit uncertainty modeling enhances predictive accuracy.



% if have a single appendix:
%\appendix[Proof of the Zonklar Equations]
% or
%\appendix  % for no appendix heading
% do not use \section anymore after \appendix, only \section*
% is possibly needed

% use appendices with more than one appendix
% then use \section to start each appendix
% you must declare a \section before using any
% \subsection or using \label (\appendices by itself
% starts a section numbered zero.)
%


%\appendices
%\section{Proof of the First Zonklar Equation}
%Appendix one text goes here.

% you can choose not to have a title for an appendix
% if you want by leaving the argument blank
%\section{hello}
%Appendix two text goes here.


% use section* for acknowledgment

\section*{Acknowledgment}

This work was supported by Grant C-SEJ-128-UGR23, funded by the Consejería de Universidad, Investigación e Innovación of the Junta de Andalucía and by the European Union through the ERDF Andalusia Program 2021-2027; project PID2023-149185OBI00 funded by MICIU/AEI/10.13039/501100011033 and by ERDF/EU. Funding for open access charge: Universidad de Granada/CBUA.


% Can use something like this to put references on a page
% by themselves when using endfloat and the captionsoff option.
\ifCLASSOPTIONcaptionsoff
  \newpage
\fi



% trigger a \newpage just before the given reference
% number - used to balance the columns on the last page
% adjust value as needed - may need to be readjusted if
% the document is modified later
%\IEEEtriggeratref{8}
% The "triggered" command can be changed if desired:
%\IEEEtriggercmd{\enlargethispage{-5in}}

% references section

% can use a bibliography generated by BibTeX as a .bbl file
% BibTeX documentation can be easily obtained at:
% http://mirror.ctan.org/biblio/bibtex/contrib/doc/
% The IEEEtran BibTeX style support page is at:
% http://www.michaelshell.org/tex/ieeetran/bibtex/

\bibliographystyle{IEEEtran}
% argument is your BibTeX string definitions and bibliography database(s)
\bibliography{IEEEabrv, refs}

%
% <OR> manually copy in the resultant .bbl file
% set second argument of \begin to the number of references
% (used to reserve space for the reference number labels box)

% biography section
% 
% If you have an EPS/PDF photo (graphicx package needed) extra braces are
% needed around the contents of the optional argument to biography to prevent
% the LaTeX parser from getting confused when it sees the complicated
% \includegraphics command within an optional argument. (You could create
% your own custom macro containing the \includegraphics command to make things
% simpler here.)
%\begin{IEEEbiography}[{\includegraphics[width=1in,height=1.25in,clip,keepaspectratio]{mshell}}]{Michael Shell}
% or if you just want to reserve a space for a photo:

%\begin{IEEEbiography}{Michael Shell}
%Biography text here.
%\end{IEEEbiography}


% You can push biographies down or up by placing
% a \vfill before or after them. The appropriate
% use of \vfill depends on what kind of text is
% on the last page and whether or not the columns
% are being equalized.

%\vfill

% Can be used to pull up biographies so that the bottom of the last one
% is flush with the other column.
%\enlargethispage{-5in}

\end{document}